% GUID v6: shared preamble of the paper and of the single-section test wrapper (testwrap.tex).
\usepackage[T1]{fontenc}
\usepackage[utf8]{inputenc}
\usepackage{lmodern}
\usepackage[margin=2.5cm]{geometry}
\usepackage{amsmath,amssymb,amsthm}
\usepackage{booktabs,tabularx,array,longtable,multirow}
\usepackage{enumitem}
\usepackage{graphicx}
\usepackage[font=small,labelfont=bf]{caption}
\usepackage[hidelinks]{hyperref}
\usepackage{url}
\usepackage{xcolor}
\usepackage[most]{tcolorbox}
\usepackage{tikz}
\usetikzlibrary{arrows.meta,positioning,shapes.geometric,fit,calc}
% Primer boxes: short background for readers from other fields, placed where the material is first needed.
% Usage: \begin{primer}{Quantum information in one page}{box:qi} ... \end{primer}
\newcounter{primer}
\newenvironment{primer}[2]{\refstepcounter{primer}\label{#2}%
  \begin{tcolorbox}[breakable,enhanced,colback=black!3,colframe=black!45,boxrule=0.5pt,arc=1pt,left=6pt,right=6pt,top=4pt,bottom=4pt,
  fonttitle=\bfseries\small,title={Primer~\theprimer: #1}]\small}{\end{tcolorbox}}
\graphicspath{{../simulation/figures/}}

\newtheorem{theorem}{Theorem}
\newtheorem{proposition}{Proposition}
\newtheorem{lemma}{Lemma}
\newtheorem{corollary}{Corollary}
\theoremstyle{definition}
\newtheorem{definition}{Definition}
\newtheorem{hypothesis}{Hypothesis}
\newtheorem*{hypothesisprime}{Hypothesis 1$'$}
\newtheorem*{hypothesisprime6}{Hypothesis 6$'$}
\theoremstyle{remark}
\newtheorem{remark}{Remark}
\newcolumntype{L}{>{\raggedright\arraybackslash}X}
\newcolumntype{P}[1]{>{\raggedright\arraybackslash}p{#1}}

\setlength{\emergencystretch}{3em}
% \chk{Q5}: check Q5 of docs/guid-v6/checks.py; \chk{gravsim:G2}: a check or result of a simulation script (see the paper's Appendix A).
\newcommand{\chk}[1]{\textsf{[#1]}}
\newcommand{\tagD}{\textsf{(D)}}
\newcommand{\tagH}{\textsf{(H)}}
\newcommand{\tagP}{\textsf{(P)}}
\newcommand{\tagS}{\textsf{(S)}}
\newcommand{\ket}[1]{\lvert #1\rangle}
\newcommand{\bra}[1]{\langle #1\rvert}
\newcommand{\Tr}{\operatorname{Tr}}
\newcommand{\dd}{\mathrm{d}}
% verdicts
\newcommand{\vd}[1]{\emph{#1}}
